\section{Introduction}

Large language models (LLMs) have exhibited astounding capabilities~\cite{qin2023chatgpt,openai2023gpt4,shu2023llasm} as versatile task-solving agents, endowed with a rich blend of knowledge and skills. Nevertheless, they still face difficulties~\cite{qin2023chatgpt,openai2023gpt4,gpt4-early-exp} in tackling various tasks that require intensive knowledge and reasoning, such as avoiding hallucination~\cite{maynez-etal-2020-faithfulness}, employing slow-thinking strategies~\cite{sloman1996empirical}, ensuring trustworthiness~\cite{wang2023decodingtrust}, and in combining diverse domain knowledge and long-horizon planning. 
In contrast, humans often exploit the benefits of collaborative problem solving, which enables them to work together effectively to solve non-routine problems in diverse domains and enhance the quality and reliability of the solutions by distributing the workload among specialties and applying a diversity of perspectives and expertise~\cite{nelson2013collaborative,roschelle1995construction,barron2000achieving}.

Inspired by collaborative problem solving, several recent works~\cite{wang2023unleashing,du2023improving,liang2023encouraging,hao2023chatllm} have improved the task-solving capabilities of LLMs by integrating multi-agent discussion. However, most of these multi-agent systems depend on handcrafted or user-specified agents, with specific roles and necessitating human supervision, which often restricts the scope of collaborative applications. Moreover, manually creating a large number of experts often consumes a lot of resources. In order to adaptively solve more complex problems, this paper aims to explore a method of adaptively generating task experts and completing different tasks through multi-level collaborative cooperation among multiple experts.

In this paper, we propose AutoAgents, an innovative framework that adaptively generates and coordinates multiple specialized agents to construct an AI team according to different tasks.
Figure~\ref{fig:framework} provides a high-level overview of AutoAgents.
By generating multiple agents with distinct expert roles, we aim to form a collaborative entity that can accomplish complex tasks by leveraging the complementary strengths of each agent.
As shown in Figure~\ref{fig:process}, the process of AutoAgents is divided into two critical stages: \textbf{Drafting Stage} and \textbf{Execution Stage}. 
The drafting stage involves a collaborative discussion among three predefined agents (\textbf{Planner}, \textbf{Agent Observer}, and \textbf{Plan Observer}) to synthesize a customized agent team and an execution plan that suit the input problem or task. The execution stage refines the plan through inter-agent collaboration and feedback, and produces the final outcome. We propose self-refinement by individual agents and collaborative refinement by multiple agents to enhance agent proficiency and promote knowledge-sharing among agents. To facilitate the specific division of labor among the agents in the synthesized team, we introduce a predefined agent (\textbf{Action Observer}) to assist the agent team in sharing information, coordinating actions, reaching consensus, and adapting to the environment.

To synthesize heterogeneous information from diverse domains is often a crucial requirement in creative industries and other real-world scenarios. We illustrate a concrete example of how AutoAgents tackles the challenging task of writing a novel about the awakening of artificial intelligence in Figure~\ref{fig:framework}. The Story Planner and Researcher collaborate to devise the plot of the story with their respective expertise, while the Character Developer and Writer enrich the novel content through imagination based on the story. Moreover, we conduct quantitative experiments and case studies in complex tasks to demonstrate the effectiveness of AutoAgents. We also conduct a comprehensive analysis and demonstrate the importance of dynamic agents for handling complex tasks, the indispensability of self-refinement for proficient agents, and the effectiveness of collaborative conversation.

\begin{figure}[!t]
\centering

\includegraphics[width=\linewidth]{./figures/framework.jpg}
\caption{A schematic diagram of AutoAgents. The system takes the user input as a starting point and generates a set of specialized agents for novel writing, along with a corresponding execution plan. The agents collaboratively carry out the tasks according to the plan and produce the final novel. Meanwhile, an observer monitors the generation and execution of the Agents and the plan, ensuring the quality and coherence of the process.
}
\label{fig:framework}
\end{figure}


To summarize, this paper makes the following novel contributions: 
 \textbf{First}, we propose AutoAgents, a novel framework that dynamically synthesizes and coordinates multiple expert agents to form customized AI teams for diverse tasks. \textbf{Second}, we conduct rigorous quantitative experiments on two challenging tasks and demonstrate that AutoAgents significantly improves both knowledge acquisition and reasoning ability in LLMs and outperforms other generated-agent frameworks. \textbf{Third}, we showcase AutoAgents’ ability to adapt to complex tasks by applying it in various scenarios such as software development. \textbf{Finally}, we conduct a thorough investigation and reveal the importance of dynamic agents for accommodating complex tasks and the necessity of self-refinement for proficient agents, and the efficacy of collaborative conversation.